% Generated draft: requires mathematical and source review.
\documentclass[11pt]{article}
\usepackage[margin=25mm]{geometry}
\usepackage{fontspec}
\setmainfont{Latin Modern Roman}
\usepackage{amsmath,amssymb,mathtools,unicode-math}
\setmathfont{Latin Modern Math}
\usepackage{enumitem,needspace,xurl,hyperref}
\setlength{\parindent}{0pt}
\setlength{\parskip}{5pt}
\setlength{\emergencystretch}{3em}
\newcommand{\N}{\mathbb{N}}
\newcommand{\Z}{\mathbb{Z}}
\newcommand{\Q}{\mathbb{Q}}
\newcommand{\R}{\mathbb{R}}
\newcommand{\C}{\mathbb{C}}
\newcommand{\F}{\mathbb{F}}
\newcommand{\A}{\mathbb{A}}
\newcommand{\PP}{\mathbb{P}}
\newcommand{\E}{\mathbb{E}}
\newcommand{\eps}{\varepsilon}
\newcommand{\dd}{\,\mathrm{d}}
\begin{document}
\section*{UnsolvedMath ID: 30001316}
\textbf{Source problem code:} catalogue entry
\subsection{Definitions and mathematical statement}
Let \(m\in\Z\) with \(m>2\). The supplied record concerns solutions of
\[
(-\Delta)^m u=\lambda u\exp(mu),
\]
where \(\Delta\) denotes the Laplace operator and \(\lambda\) is a parameter. The ambient dimension, domain, boundary conditions, solution class, sign restrictions, and permitted values or limiting behavior of \(\lambda\) are not specified.

The record refers to solutions with ``bounded concentration energy,'' a ``concentration mass,'' and a ``fundamental level'' \(\Lambda_1\), without defining these quantities or the concentrating family involved. If \(M\) denotes the intended concentration mass, its proposed quantization conclusion is
\[
\exists L\in\Z_{>0}\qquad M=L\Lambda_1.
\]
The question is whether this conclusion holds for every concentrating family satisfying the intended energy bound when \(m>2\). This is a schematic statement: the supplied information does not determine the energy functional, the concentration procedure, the normalization of \(\Lambda_1\), or whether \(M\) is the mass at an individual concentration point or a total concentration mass. It also does not specify whether only nonzero concentration masses are under consideration, as the requirement \(L>0\) would suggest.

The equation above retains the supplied exponent \(mu\) exactly. The designation ``critical'' cannot be made precise without additional dimensional and functional hypotheses. The supplied record states that the analogous quantization holds for \(m=1,2\), but gives no precise hypotheses for that comparison. Consequently, a fully quantified mathematical formulation requires consultation of the cited source.
\subsection{Short English statement}
For the supplied higher-order equation, the question asks whether concentration masses under an unspecified energy bound are positive integer multiples of a fundamental level when \(m>2\). The definitions of energy, concentration mass, fundamental level, and the analytic setting are missing, so this remains a statement-only draft requiring mathematical review.
\subsection{Sources}
\begin{itemize}
\item Local UnsolvedMath registry: permanent ID \(30001316\), exact source problem code \texttt{catalogue entry}, title ``Quantization of Concentration Mass for Higher-Order Equations,'' in the Oberwolfach Reports Open Problems collection. No individual proposer is supplied; the proposed year is recorded as 2009.
\item Supplied source citation: \emph{Partielle Differentialgleichungen} (2010), \emph{Oberwolfach Reports}, pp.\ 1973--2024, DOI 10.4171/owr/2009/36. \url{https://doi.org/10.4171/owr/2009/36}. The supplied literature-source label is ``original Oberwolfach report 4079.'' The background describes the reports as records of mathematical workshops and their open-problem discussions.
\item Supplied literature source labeled ``Martinazzi threshold paper'': \url{https://cvgmt.sns.it/paper/1547/}. No title, full author list, publication details, or precise theorem statement is supplied for this reference.
\item Supplied literature assessment, checked on exactly 2026-08-21: the record labels the problem open and reports that the cited higher-order work provides threshold results without a general resolution of this quantization question. This is supplied dated curation derived from the OWR material, not an independently verified literature assessment; its status and cited claims require independent verification.
\item The registry labels the statement \texttt{corrected\_verified}; its supplied verification explanation says that the complete statement was preserved while mathematical notation was placed in delimiters. This label does not supply the missing analytic hypotheses or establish mathematical correctness.
\end{itemize}
\end{document}
